\section{敏感性分析}

\subsection{储能效率口径}

题目中的“充放电效率 90\%”可能解释为单向效率，也可能解释为往返效率。本文以单向效率 0.9 为主值，并在问题一中改变效率，结果见表~\ref{tab:eta}。

\begin{table}[H]
\centering
\caption{问题一对单向效率的敏感性}
\label{tab:eta}
\begin{tabular}{rrrr}
\toprule
$\eta$ & 购电费/元 & 充电量/kWh & 放电量/kWh \\
\midrule
0.850000 & 36603.41 & 22295.51 & 16108.50 \\
0.900000 & 35126.95 & 20740.67 & 16799.94 \\
$\sqrt{0.9}=0.948683$ & 33801.50 & 19842.71 & 17858.44 \\
0.950000 & 33767.03 & 19801.75 & 17871.08 \\
\bottomrule
\end{tabular}
\end{table}

效率越高，单位外购电可转移至高价时段的有效电量越多，费用单调下降。从 $0.9$ 改为
$\sqrt{0.9}$ 会使问题一费用下降 $1325.45$ 元，表明效率口径会显著影响绝对结果，故必须在提交时明确。

\subsection{安全分位数}

1 月滚动验证中，安全分位数从 0.70 提高到 0.75 时加权损失由 234008.90 降至
229079.15；继续提高至 0.80、0.85 和 0.90 时，损失分别增至 233219.96、
248703.67 和 279456.50。结果呈先降后升趋势：裕度过小会产生 5 倍紧急费用，裕度过大则造成计划过购和弃电。0.75 是给定数据下的验证最优值，并非人为固定参数。

\subsection{终端规则与价格信息}

问题二至四采用跨日连续 SOC，而非每日强制回到初值。最后一日前计划虽强制目标
$6000~\mathrm{kWh}$，实际执行仍受预测误差影响，四个主场景年末 SOC 分别为
6208.71、6635.92、6221.18 和 6644.19 kWh。本文将此偏差作为终端风险披露，不利用未来实际值事后修饰计划。

波动电价下，完美价格信息比因果价格预测节省 5.13 万元，说明价格信息改进的收益有限但非零。相比之下，是否增加日内光伏预报对全年费用的影响达到数十万元，是更关键的策略参数。
